\setlength{\absparsep}{18pt} % ajusta o espaçamento dos parágrafos do resumo
\begin{resumo}
O monitoramento da biodiversidade costeira é essencial para a conservação ambiental, porém frequentemente enfrenta desafios relacionados à escala de coleta e à qualidade dos dados. Este trabalho, em colaboração com o Centro de Estudos Costeiros Limnológicos e Marinhos (CECLIMAR), teve como objetivo desenvolver uma aplicação móvel, denominada \textit{FaunaMar}, visando automatizar, padronizar e otimizar o ciclo de vida dos dados de monitoramento da fauna costeira no litoral do Rio Grande do Sul.

A metodologia utilizada foi o modelo de desenvolvimento iterativo incremental com práticas ágeis, gerenciado com Kanban pela plataforma \textit{Jira}. Foram realizadas reuniões com profissionais do CECLIMAR para levantamento de requisitos e a prototipação foi conduzida na ferramenta \textit{Figma}. A aplicação foi implementada em \textit{Flutter} com backend em \textit{Firebase}, incluindo autenticação, banco de dados NoSQL, armazenamento de imagens e distribuição de versões de teste. 

Os resultados incluem uma aplicação funcional que digitaliza e centraliza todo o processo de monitoramento. O sistema possui dois perfis de usuário: o cientista cidadão e o pesquisador, com interfaces e permissões distintas. Entre as principais funcionalidades estão: formulários de registro simples e técnico, captura de geolocalização por GPS e painel de registros com filtros e mapa interativos. Há também um perfil gamificado com conquistas, além de um módulo exclusivo para que os pesquisadores avaliem os registros. O projeto teve um total de 131 tarefas, das quais 86\% foram concluídas, demonstrando um processo de desenvolvimento estruturado e rastreável.

A aplicação foi testada em prática por usuários, validando sua usabilidade e viabilidade. Entregando uma solução tecnológica que atende a uma demanda real do CECLIMAR com a aplicação da Ciência Cidadã como ferramenta para a pesquisa ambiental. O aplicativo FaunaMar auxilia na qualidade e na eficiência do monitoramento da fauna costeira, reduzindo a carga operacional dos pesquisadores e fortalecendo o vínculo entre a ciência e a sociedade para a promoção de um futuro mais sustentável.

\textbf{Palavras-chave}: ciência cidadã. monitoramento ambiental. aplicativo móvel.
\end{resumo}